\begin{abstract}
针对无人机烟幕干扰弹在来袭导弹与真目标之间实施视线遮蔽的问题，本文建立统一的三维运动学、有限视线段几何判定和多平台时空协同优化模型。与仅用目标中心点判断遮蔽的方法不同，本文保留半径 $7\,\mathrm{m}$、高 $10\,\mathrm{m}$ 的完整圆柱目标，以烟幕球心到导弹--圆柱各条视线段的最大距离不超过 $10\,\mathrm{m}$ 作为完全遮蔽条件，并通过事件求根计算有效时间集合。

对问题一，按照给定参数计算单枚烟幕弹的完全遮蔽时长，并完成边界离散与重力参数检验。对问题二，将航向、速度、投放时刻和引信延时作为连续变量，采用多初值差分进化与局部细化求解，得到 FY1 单弹最大有效时长为 $4.588062\,\mathrm{s}$。对问题三，建立共享航迹和最小投放间隔约束，以三枚烟幕弹有效区间并集为目标，通过错峰起爆把三个窗口拼接为 $[5.205021,12.814687]$，联合时长为 $7.609665\,\mathrm{s}$。

对问题四，将三机协同分解为单弹窗口搜索与联合时间协调。FY1、FY2、FY3 分别形成早、中、晚三个互不重叠的窗口，M1 的联合遮蔽时长达到 $10.657048\,\mathrm{s}$。对问题五，构造含二进制任务分配、单机容量、共享航迹、投放间隔、起爆高度和烟幕寿命的混合优化模型；先进行无人机--导弹效能探针，再按边际收益分配并分组连续优化。最终使用 $8$ 枚有效烟幕弹，M1、M2、M3 的联合时长分别为 $7.609665$、$10.647375$ 和 $3.721573\,\mathrm{s}$，总目标为 $21.978613\,\mathrm{s}$。

终算采用每个圆周 $5760$ 个角向点、$0.005\,\mathrm{s}$ 扫描与 Brent 事件求根，全部策略均通过速度、投放间隔、起爆高度和导弹到达时限回代。模型能够统一处理单弹、多弹、多机和多目标情形，并为带时空窗口的多平台遮蔽与覆盖任务提供可复用框架。

\keywords{烟幕干扰\quad 视线遮蔽\quad 区间并集\quad 差分进化\quad 任务分配}
\end{abstract}
\endinput
\begin{abstract}
\keywords{烟幕干扰\quad 视线遮蔽\quad 弹道模型\quad 事件求根}
\end{abstract}
\endinput

\begin{abstract}
\textit{cumcmthesis} 是为\textbf{中国大学生数学建模竞赛}(\textbf{CUMCM, China Undergraduate Mathematical Contest in Modeling})编写的 \LaTeX{} 模板, 旨在让大家专注于论文的内容写作, 而不用花费过多精力在格式的定制和调整上.  使用者需要有一 定的 \LaTeX{} 的使用经验, 至少要会使用常用宏包的一些功能, 比如参考文献，数学公式，图片使用，列表环境等等.

% 关键词
\keywords{\TeX{}\quad \LaTeX{}\quad 图片\quad 表格\quad 公式}
\end{abstract}
